\documentclass[10pt,a4paper]{amsart}
\usepackage[margin=19mm]{geometry}
\usepackage{amsmath,amssymb,mathtools}
\usepackage[scr=rsfs,cal=euler]{mathalfa}
\usepackage{xfrac}
\usepackage[unicode,hidelinks,bookmarks=false]{hyperref}
\newcommand{\ie}{i.e.\ }
\newcommand{\cf}{cf.\ }
\newcommand{\resp}{respectively\ }
\newcommand{\Subsection}{\subsection}
\providecommand{\nobreakdash}{\nobreak\hbox{-}}
\setlength{\emergencystretch}{2em}
\multlinegap0pt
\usepackage{amssymb}
\usepackage{mathtools}
\usepackage{bm}
\usepackage{tikz}
\usepackage{xcolor}
\newtheorem{proposition}{Proposition}
\newcommand{\R}{\mathbb R}
\newcommand{\diag}{\mathrm{diag}}
\newcommand{\rh}[1]{\textcolor{purple}{\textsf{[RH: #1]}}}
\title{\LARGE \bf Information Revelation and Alignment Faking in Stochastic Differential Games}
\author{Daniel Ralston, Xu Yang, Ruimeng Hu}
\date{Source: arXiv:2603.17197v2 (CC BY 4.0)}
\begin{document}
\maketitle

\noindent\emph{Source and licence.} \LARGE \bf Information Revelation and Alignment Faking in Stochastic Differential Games. Version arXiv:2603.17197v2 (CC BY 4.0), distributed by arXiv under the Creative Commons Attribution 4.0 International licence, http://creativecommons.org/licenses/by/4.0/, as recorded in the arXiv licence metadata for this exact version and shown on the abstract page https://arxiv.org/abs/2603.17197v2. Full licence notice: \texttt{licenses/arxiv-2603.17197-CC-BY-4.0.txt}.\par\smallskip

\noindent\textit{Editorial scope.} Counted unit: the article's proposition prop:exist\_unique\_baseline (source0.tex lines 321--324). The article's complete original proof of it is delivered with it (source0.tex lines 326--332, one \texttt{proof} environment of 7 lines). No mathematics is written, repaired or re-derived: the statement and the proof are the article's own lines, in the article's own notation, and no retained line is substituted or re-flowed.  Retained uncounted as the source-local context the proof needs: lines 269--273; lines 290--299; lines 292--295.  Nothing was omitted from any retained span, so the package carries no omission record.  No retained line contains a citation, so the wrapper carries no bibliography and no bibliography entry.  No retained line contains a figure environment, an image-inclusion command or drawing code, so no image is delivered and the package declares no asset dependency.  Editorial formatting only, in addition: an AMS \texttt{amsart} preamble that restates the article's own declarations and macros used by the retained lines, namely the environments \texttt{proposition} on separate counters, together with the standard packages those lines need.  The retained environments print in this document as Proposition 1 in order of appearance, whereas the article prints the same items with the numbering of its own full document.  The complete native source of the article as distributed in the arXiv e-print for this exact version is bundled unchanged as \texttt{sources/196609/source0.tex}.
\begin{align}\label{eq:Riccati_A}
    \dot{\theta}^A(t) &= \theta^A R_A \theta^A + \theta^A R_B \theta^B + \theta^B R_B \theta^A - Q_A, \\
    \label{eq:Riccati_B}
    \dot{\theta}^B(t) &= \theta^B R_B \theta^B + \theta^A R_A \theta^B + \theta^B R_A \theta^A - Q_B,
\end{align}

\begin{proposition}\label{prop:theta_bound}
    Let $\theta = \diag\{\theta^A, \theta^B\}$, and suppose \vspace{-0.3em}    
    \begin{multline}\label{eq:T_bound}
        T \leq \mathcal{T}(m_A, m_B)
        := (c_1 c_2)^{-1/2}  \\ \cdot\arctan \big((1+m_A^2 + m_B^2)(c_2/c_1)^{-1/2}\big),
    \end{multline}
    where $c_1 = 5 (r_A^{-2} + r_B^{-2})^{1/2}$ and $c_2= \sqrt{q_A^2 + q_B^2}\,(1 + m_A^2 + m_B^2).$ 
    Then any solution to \eqref{eq:Riccati_A}--\eqref{eq:Riccati_B} on $[0,T]$ satisfies $\big\|\theta(t; m_A, m_B)\big\|_F \le 1 + m_A^2 + m_B^2$. 
   % Furthermore, assume a solution of $\theta$ exists for $0\leq t \leq T$.  Then, 
\end{proposition}

    \begin{multline}
        T \leq \mathcal{T}(m_A, m_B)
        := (c_1 c_2)^{-1/2}  \\ \cdot\arctan \big((1+m_A^2 + m_B^2)(c_2/c_1)^{-1/2}\big),
    \end{multline}

\begin{proposition}\label{prop:exist_unique_baseline}
    Let $T \leq \mathcal{T}(m_A, m_B)$ as in \eqref{eq:T_bound}.  Then the coupled Riccati system \eqref{eq:Riccati_A}--\eqref{eq:Riccati_B} for $\theta^A$, $\theta^B$ %with terminal conditions $\theta^A(T) = \theta^B(T) = 0$ 
    admits a unique solution over $[0,T]$.
\end{proposition}

\begin{proof}
    Using the same definitions as in Proposition~\ref{prop:theta_bound}, let $Y(t) = \theta(T-t)$ so that $Y$ solves the initial value problem $\dot{Y} = -F(Y) + Q$, $Y(0) = 0$.  
    The right-hand side $G(Y) := -F(Y) + Q$ is quadratic in $Y$, and is therefore locally Lipschitz on any bounded subset of $\R^{4\times 4}$.
    By the Picard--Lindel\"of theorem, there exists a unique solution of $Y$ on $[0,T^\ast)$, where $T^\ast$ is maximal.
    Suppose that $T^* \leq T$; then $\|Y(t)\|_F \to \infty$ as $t \to T^*$. 
    However, on $[0, T^\ast)$ where the solution of $Y(t)$ exists, Proposition~\ref{prop:theta_bound} applies, yielding $\|Y(t)\|_F \leq 1 + m_A^2 + m_B^2$ for all $t \in [0,T^\ast - \epsilon]$, contradicting blowup at $T^*$. %\rh{I change the wording to "By the Picard--Lindelhof theorem, there exists a unique solution of $Y$ on $[0,T^\ast)$, where $T^\ast$ is maximal."}
\end{proof}

\end{document}
